\documentclass[11pt,a4paper]{article}
\usepackage[T1]{fontenc}
\usepackage[utf8]{inputenc}
\usepackage[english]{babel}
\usepackage{lmodern,amsmath,amssymb,mathtools,bm}
\usepackage[margin=2.2cm]{geometry}
\usepackage{booktabs,longtable,array,listings,xcolor}
\usepackage[hidelinks]{hyperref}
\newcommand{\ket}[1]{\left|#1\right\rangle}
\newcommand{\bra}[1]{\left\langle#1\right|}
\newcommand{\braket}[2]{\left\langle#1\middle|#2\right\rangle}
\newcommand{\Tr}{\operatorname{Tr}}
\newcommand{\ii}{\mathrm{i}}
\newcommand{\code}[1]{\texttt{\detokenize{#1}}}
\newcommand{\B}{\mathcal B}
\newcommand{\K}{\mathcal K}
\lstset{language=Python,basicstyle=\ttfamily\scriptsize,numbers=left,
 numberstyle=\tiny\color{gray},breaklines=true,columns=fullflexible,
 keepspaces=true,showstringspaces=false,frame=single,tabsize=4}
\setlength{\parindent}{0pt}
\setlength{\parskip}{5pt}
\setlength{\emergencystretch}{2em}
\setcounter{tocdepth}{2}
\title{Multimode Rabi dynamics:\\finite models, experiments and patch analysis}
\author{Documentation of the \texttt{rabi-truncation} repository}
\date{28 September 2026}
\begin{document}
\maketitle
\begin{abstract}
The repository evolves pure states with one shared engine, enumerated photon
occupation bases and sparse Hamiltonians. This document specifies the finite
models actually solved, derives their preparation and measurement formulas,
explains the common-space embedding used for fidelities, and documents the
Python and notebook interfaces. The comparison with the student implementation
separates unitary changes of convention, changes of the finite physical model,
and corrections to numerical comparisons. Numbered source listings are followed
by comments tied to their current lines. The common \code{full} space used for
overlaps is implicit; choosing a separately propagated comparison trajectory
does not require propagating or allocating that ambient space.
\end{abstract}
\tableofcontents

\section{Finite model, state preparation and calculated quantities}
\subsection{Periodic grid, parameter controls and photon spaces}
\paragraph{Notation and units.}
Throughout, $\hbar=c=1$. Length and time have inverse-energy units.
The TLS energy origin is its ground state: its free Hamiltonian is
$\omega_0\ket e\bra e$, rather than $\omega_0\sigma_z/2$.
These differ by a scalar, which matters when reporting an absolute energy.
The symbols used below are:
\begin{longtable}{p{.18\textwidth}p{.73\textwidth}}
\toprule Symbol & Meaning in the implementation\\\midrule\endhead
$L$ & Periodic box length, supplied in \code{param_atom}.\\
$k_m,\omega_m$ & Signed momentum and nonnegative frequency $\omega_m=|k_m|$.\\
$M$ & Number of retained modes; after selection this means the selected
count, unless explicitly described as the base-grid count.\\
$N,B,d$ & \code{n_max}, number of photon configurations, and joint ket
dimension $d=2B$.\\
$\bm n,\nu$ & Occupation tuple and total photon number $\nu=\sum_m n_m$.\\
$D,f_m,u_m$ & Physical coupling, chosen profile and complex annihilation
coefficient, defined in equation~\eqref{eq:H}.\\
$k_0,\sigma_k,x_0$ & Packet momentum centre, width convention and initial
Fourier position parameter.\\
$x_{\mathrm{tls}}$ & Parameter in the imposed annihilation phase;
its plotted interaction coordinate is derived in equation~\eqref{eq:position}.\\
$T,\code{dt},N_t$ & Final time, nominal output spacing and number of
output times, including both endpoints.\\
\bottomrule
\end{longtable}
\paragraph{Grid and the two controls.}
For a periodic box of length $L>0$,
\begin{equation}
 \varphi_j(x)=\frac{e^{\ii k_jx}}{\sqrt L},\qquad
 \Delta k=\frac{2\pi}{L},\qquad k_j=j\Delta k,\quad j\in\mathbb Z,
 \qquad \omega_j=|k_j|.
 \label{eq:grid}
\end{equation}
Indeed, $\varphi_j(x+L)=\varphi_j(x)$ and
$\int_{x_*}^{x_*+L}\varphi_j^*\varphi_\ell\,dx=\delta_{j\ell}$.
Signed momenta enter spatial phases; their absolute values enter the dispersion
and square-root profile. Modes are sorted by signed momentum. There is no
automatic deletion of the zero mode.

When \code{CTRL_M_EXPLICIT=True}, the inputs are $L$ and a positive odd
integer $M$. Setting $J=(M-1)/2$ gives
\begin{equation}
 \K_M=\{j\Delta k:-J\leq j\leq J\},\quad |\K_M|=M,
 \quad k_{\mathrm{IR,eff}}=0,\quad
 k_{\mathrm{UV,eff}}=\frac{M-1}{2}\Delta k.
 \label{eq:explicit}
\end{equation}
The cutoff dictionary is inactive. An even $M$ raises \code{ValueError};
it is never rounded, decremented or made asymmetric. $M=1$ gives
$\K=\{0\}$, not a resonant single-mode cavity at $\omega_0$.
Booleans and floating-point values, including \code{5.0}, are rejected as
mode-count inputs; Python and NumPy integer types are accepted.

When \code{CTRL_M_EXPLICIT=False}, the inputs are $L$ and finite radial
cutoffs $0\leq k_{\mathrm{IR}}\leq k_{\mathrm{UV}}$:
\begin{equation}
 \K_{\mathrm{IR,UV}}=\{j\Delta k:
 k_{\mathrm{IR}}\leq |j|\Delta k\leq k_{\mathrm{UV}}\}.
 \label{eq:cutoffs}
\end{equation}
For strictly positive indices, in exact arithmetic,
\begin{align}
 j_{\min}&=\max\left(1,\left\lceil\frac{k_{\mathrm{IR}}}{\Delta k}\right\rceil\right),&
 j_{\max}&=\left\lfloor\frac{k_{\mathrm{UV}}}{\Delta k}\right\rfloor,\\
 M_{\mathrm{eff}}&=2\max(0,j_{\max}-j_{\min}+1)
                       +\mathbf1_{k_{\mathrm{IR}}=0}.
 \label{eq:count}
\end{align}
The code subtracts $10^{-12}$ before the ceiling and adds $10^{-12}$ before
the floor, protecting numerically boundary-aligned points. This tolerance is
on the dimensionless cutoff-to-spacing ratios. Zero is included if and only
if the input IR equals zero exactly. A positive IR excludes it through the
band definition even when that IR is very small. An empty band is an error.
The resulting count can be even when zero is excluded; the odd-$M$ condition
belongs only to the explicit-count branch.

Thus specifying $L$ and $M$ implicitly fixes a UV cutoff, while specifying
$L$ and cutoffs implicitly fixes the count. These descriptions coincide for
a symmetric contiguous grid containing zero, but not for arbitrary selected
subsets or positive-IR bands. Continuous cutoffs are not uniquely recoverable
from a mode set: with IR zero, the same $\K_M$ is obtained whenever
\begin{equation}
 J\Delta k\leq k_{\mathrm{UV}}<(J+1)\Delta k
 \label{eq:uvinterval}
\end{equation}
in exact arithmetic. The estimator reports effective radial extrema, without
claiming that the requested cutoff originally equalled an extremum. It reports
the smallest nonzero momentum separately: effective IR zero does not mean
zero spacing between the nonzero modes.

\paragraph{Additional mode selection.}
Selection is applied after the base grid is resolved. With
$w_p=\code{photon_window}$ and $w_a=\code{atom_window}$ it retains the union
\begin{equation}
 |k-k_0|\leq w_p\sigma_k\quad\text{or}\quad
 |k-\omega_0|\leq w_a\sigma_k\quad\text{or}\quad
 |k+\omega_0|\leq w_a\sigma_k.
 \label{eq:selection}
\end{equation}
Each comparison has an absolute $10^{-12}$ momentum tolerance. Each
individually empty window adds its nearest available mode; ties select the
first mode in the sorted array. A zero-radius window therefore still retains
a nearest mode without an exact resonance. Factors must be finite and
nonnegative. Selection can yield asymmetry, holes, an even selected count,
or no zero mode. These follow from the requested windows, not a hidden
modification of the base-grid convention.

The packet is then rebuilt and normalized on the selected grid. A comparison
therefore tests both preparation and dynamics. For a selected subset
$S\subseteq\K$, with $c_m$ the normalized unselected packet,
\begin{equation}
 p_S=\sum_{m\in S}|c_m|^2,\qquad
 c_m^{(S)}=c_m/\sqrt{p_S}\ (m\in S),\qquad
 F_{\mathrm{state}}(0)=p_S,\quad F_{\mathrm{atom}}(0)=1
 \label{eq:selectedprep}
\end{equation}
for identical number preparations with $1\leq n\leq N$, defined below.
For a coherent preparation, renormalizing the selected packet also changes
its individual mode amplitudes. The fixed-product projection formula below
does not automatically apply across such different grids.

\paragraph{Three photon bases and their ordering.}
A photon configuration is $\bm n\in\mathbb N^M$. With $N=\code{n_max}$,
\begin{align}
 \B_{\mathrm{truncated}}&=\{\bm0\}\cup
        \{n\bm e_m:1\leq m\leq M,\ 1\leq n\leq N\},\\
 \B_{\mathrm{total}}&=\{\bm n:\textstyle\sum_m n_m\leq N\},\\
 \B_{\mathrm{full}}&=\{\bm n:0\leq n_m\leq N\ \forall m\}.
 \label{eq:bases}
\end{align}
The string names are \code{truncated}, \code{full+totalcap} and \code{full}.
The TLS is a separate $\mathbb C^2$ factor in all three. There is no supplemental
atomic oscillator or \code{truncated+atom} basis. At fixed grid and cap,
\begin{equation}
 \B_{\mathrm{truncated}}\subseteq\B_{\mathrm{total}}\subseteq\B_{\mathrm{full}}.
\end{equation}
The joint ket dimensions are
\begin{align}
 d_{\mathrm{truncated}}&=2(1+MN),&
 d_{\mathrm{total}}&=2\binom{M+N}{N},&
 d_{\mathrm{full}}&=2(N+1)^M.
 \label{eq:dimensions}
\end{align}
Configurations with total $r$ are weak compositions of $r$ into $M$ entries,
numbering $\binom{M+r-1}{r}$. Summing over $0\leq r\leq N$ gives
$\binom{M+N}{N}$. At $N=0$ every basis has one vacuum and two TLS components;
at $N=1$ truncated and total-cap coincide. For $M=1$ all three coincide.

Indices are representation dependent. For $M=2,N=2$, the exact source order is:
\begin{center}
\begin{tabular}{c c c c}
\toprule $i$ & \code{truncated} & \code{full+totalcap} & \code{full}\\\midrule
0 & $(0,0)$ & $(0,0)$ & $(0,0)$\\
1 & $(1,0)$ & $(1,0)$ & $(0,1)$\\
2 & $(2,0)$ & $(0,1)$ & $(0,2)$\\
3 & $(0,1)$ & $(2,0)$ & $(1,0)$\\
4 & $(0,2)$ & $(1,1)$ & $(1,1)$\\
5 & --- & $(0,2)$ & $(1,2)$\\
6 & --- & --- & $(2,0)$\\
7 & --- & --- & $(2,1)$\\
8 & --- & --- & $(2,2)$\\\bottomrule
\end{tabular}
\end{center}
Total-cap enumerates combinations of mode indices with replacement and converts
their multiplicities to occupations. Full uses a Cartesian product with the
last entry changing fastest. In every representation, $s=0$ denotes $g$,
$s=1$ denotes $e$, and
\begin{equation}
 q=2i+s,\qquad v_{2i}=C_{i,g},\quad v_{2i+1}=C_{i,e},\qquad
 \ket\psi=\sum_{i,s}C_{i,s}\ket{\bm n_i;s}.
 \label{eq:index}
\end{equation}
An occupation-to-index dictionary supplies targets. Equal raw indices in
different bases do not generally mean equal physical configurations.

\subsection{Preparation, projected Hamiltonian and evolution}
\paragraph{Gaussian coefficients and number preparation.}
For $k_0,x_0\in\mathbb R$ and $\sigma_k>0$,
\begin{equation}
 \widetilde c_m=\exp\left[-\frac{(k_m-k_0)^2}{4\sigma_k^2}\right]
                  e^{-\ii k_mx_0},\qquad
 c_m=\frac{\widetilde c_m}{\sqrt{\sum_\ell|\widetilde c_\ell|^2}}.
 \label{eq:packet}
\end{equation}
All selected signed modes enter. There is no right-moving-only option:
positive $k_0$ specifies a centre, not a filter. In the continuous ideal
Gaussian, $|\widetilde c(k)|^2$ has standard deviation $\sigma_k$.
The actual finite-band moments are
\begin{equation}
 \langle k\rangle_c=\sum_m k_m|c_m|^2,\qquad
 (\Delta k_c)^2=\sum_m(k_m-\langle k\rangle_c)^2|c_m|^2,
\end{equation}
and need not equal $k_0$ and $\sigma_k^2$. The code subtracts the maximum
real exponent before exponentiation. This common positive scale cancels at
normalization and avoids an entirely underflowed packet for a distant centre.

The TLS preparation is one of
\begin{equation}
 \ket{a_0}\in\left\{\ket g,\ket e,
 (\ket g+\ket e)/\sqrt2,(\ket g-\ket e)/\sqrt2\right\},
\end{equation}
encoded by \code{initial_state='g'}, \code{'e'}, \code{'+'} or \code{'-'}
in the atomic dictionary, with default $g$. For \code{state='number'},
\begin{equation}
 \ket{\psi_0^{(n)}}=
 \begin{cases}
 \ket{\bm0}\otimes\ket{a_0},&n=0,\\
 \sum_m c_m\ket{n\bm e_m}\otimes\ket{a_0},&1\leq n\leq N.
 \end{cases}
 \label{eq:number}
\end{equation}
The default is this number preparation with $n=1$. Orthogonality of the
configurations proves normalization for $n>0$. At $n=0$ the vacuum is
constructed once; summing it once per mode would duplicate one physical state.

For $n>1$, all $n$ photons occupy one mode at a time. A Fock state of a
collective packet mode $b^\dagger=\sum_m c_ma_m^\dagger$ would instead be
\begin{equation}
 \frac{(b^\dagger)^n}{\sqrt{n!}}\ket0
 =\sum_{\sum_m n_m=n}
       \sqrt{\frac{n!}{\prod_m n_m!}}\prod_m c_m^{n_m}\ket{\bm n}.
 \label{eq:collective}
\end{equation}
For two modes and $n=2$, this has amplitudes
$c_1^2,\sqrt2c_1c_2,c_2^2$ on $(2,0),(1,1),(0,2)$; the implementation has
$c_1,c_2$ on $(2,0),(0,2)$ only. The historical preparation is preserved,
so higher-$n$ results must use its actual definition.

\paragraph{Coherent projection and its retained weight.}
For \code{state='coherent'}, $\alpha$ is finite and complex, and
\begin{equation}
 \ket{\boldsymbol\alpha}=\bigotimes_m\ket{\alpha c_m},\qquad
 b_{\bm n}=e^{-|\alpha|^2/2}
          \prod_m\frac{(\alpha c_m)^{n_m}}{\sqrt{n_m!}}.
 \label{eq:coherentcoeff}
\end{equation}
The total unprojected mean photon number is $\lambda=|\alpha|^2$, with
individual means $\lambda_m=\lambda|c_m|^2$ and $\sum_m\lambda_m=\lambda$.
Let
\begin{equation}
 P_\B=\sum_{\bm n\in\B,s}\ket{\bm n;s}\bra{\bm n;s},\qquad
 Z_\B=\sum_{\bm n\in\B}|b_{\bm n}|^2.
\end{equation}
The actual prepared ket is
\begin{equation}
 \ket{\psi_{0,\B}}=
 \frac{P_\B(\ket{\boldsymbol\alpha}\otimes\ket{a_0})}{\sqrt{Z_\B}}.
 \label{eq:coherent}
\end{equation}
The common exponential prefactor is not evaluated in the source because it
cancels on normalization. With $E_N(z)=\sum_{r=0}^N z^r/r!$,
\begin{align}
 Z_{\mathrm{full}}&=e^{-\lambda}\prod_m E_N(\lambda_m),\\
 Z_{\mathrm{total}}&=e^{-\lambda}E_N(\lambda),\\
 Z_{\mathrm{truncated}}&=e^{-\lambda}
       \left[1+\sum_m\big(E_N(\lambda_m)-1\big)\right].
 \label{eq:weights}
\end{align}
The total-cap formula follows from the multinomial identity
\begin{equation}
 \sum_{\sum_m n_m=r}\prod_m\frac{\lambda_m^{n_m}}{n_m!}
       =\frac{(\sum_m\lambda_m)^r}{r!}=\frac{\lambda^r}{r!}.
\end{equation}
Truncated counts vacuum once and nonzero occupations on each mode ray.
Consequently $0<Z_{\mathrm{truncated}}\leq Z_{\mathrm{total}}
\leq Z_{\mathrm{full}}\leq1$. At $\alpha=0$ all weights equal one.
If several $\lambda_m$ are nonzero, increasing $N$ in truncated alone
approaches $e^{-\lambda}[1+\sum_m(e^{\lambda_m}-1)]$, generally below one:
multimode occupations remain omitted at every cap.

For two projections of the \emph{same} coherent product on the same grid,
\begin{equation}
 F_{\mathrm{state}}(0)=\frac{Z_{A\cap B}^{\,2}}{Z_AZ_B},\qquad
 Z_{A\cap B}=\sum_{\bm n\in\B_A\cap\B_B}|b_{\bm n}|^2,
 \qquad F_{\mathrm{atom}}(0)=1.
 \label{eq:initialf}
\end{equation}
If $\B_A\subseteq\B_B$, this is $Z_A/Z_B$. For two equally weighted modes,
$\lambda=1$ and $N=2$, the weights are $2.25/e$, $2.5/e$, and
$2.640625/e$ for truncated, total-cap and full. Truncated versus total-cap
therefore starts at $F_{\mathrm{state}}(0)=0.9$ before any interaction.
All these preparations are field--TLS products, explaining their unit
atomic fidelity and zero initial atomic entropy.

\paragraph{Hamiltonian and matrix entries.}
For finite real $D$, $\omega_0\geq0$ and finite $x_{\mathrm{tls}}$,
\begin{align}
 H_0&=\sum_m |k_m|a_m^\dagger a_m+\omega_0\ket e\bra e,\\
 u_m&=\frac{\ii f_m e^{-\ii k_mx_{\mathrm{tls}}}}{\sqrt L},\qquad
 f_m=\begin{cases}\sqrt{|k_m|},&\code{sqrt},\\1,&\code{flat},\end{cases}\\
 H_\B&=P_\B\left[H_0+D\sum_m(u_ma_m+u_m^*a_m^\dagger)
                         (\sigma_++\sigma_-)\right]P_\B
       \equiv H_{0,\B}+DV_\B,
 \label{eq:H}
\end{align}
where $\sigma_+=\ket e\bra g$ and $\sigma_-=\ket g\bra e$.
The factor $\ii$ remains at zero TLS position. Creation receives the conjugate
annihilation coefficient. Because field and TLS operators commute,
$(u_ma_m\sigma_x)^\dagger=u_m^*a_m^\dagger\sigma_x$; the sum is Hermitian,
and orthogonal projection preserves this property. In energy units, $D$
is dimensionless for \code{sqrt} and has units $\sqrt{\text{energy}}$ for
\code{flat}. Equal numerical couplings under the profiles do not describe
the same physical model.

For a source column $q=2i+s$, the assembled entries are
\begin{align}
 (H_0)_{2i+s,2i+s}&=\sum_m|k_m|n_{i,m}+s\omega_0,\\
 (V_\B)_{2j+1-s,2i+s}^{(-)}&=u_m\sqrt{n_{i,m}},
       &&\bm n_j=\bm n_i-\bm e_m,\\
 (V_\B)_{2j+1-s,2i+s}^{(+)}&=u_m^*\sqrt{n_{i,m}+1},
       &&\bm n_j=\bm n_i+\bm e_m\in\B.
 \label{eq:entries}
\end{align}
A descending transition needs $n_{i,m}>0$. All bases are downward closed,
so the target then exists. An ascending target is looked up and omitted if
absent, realizing exactly $P_\B H P_\B$. No omitted target is replaced by
another oscillator. A transition and its reverse are both retained or both
absent, with conjugate coefficients.

For example, creation in mode two from $(1,0)$ targets $(1,1)$: retained
in total-cap/full at $N=2$, absent in truncated. Creation from $(2,0)$ into
$(2,1)$ is retained only in full. Shared entries have identical coefficients,
but different transition graphs can produce different dynamics.
Projected ladder operators are not canonical at the boundary. For one mode,
\begin{equation}
 a_N\ket n=\sqrt n\ket{n-1},\quad a_N^\dagger\ket N=0,\qquad
 [a_N,a_N^\dagger]=I-(N+1)\ket N\bra N.
 \label{eq:boundary}
\end{equation}
A calculation can be exact for this finite model without proving convergence
to unbounded Fock space. Assembly costs $O(Md)$, rather than scanning all pairs.

\paragraph{RWA, parity and the zero mode.}
Under RWA only annihilation with excitation and creation with de-excitation remain:
\begin{equation}
 H_{\B,\mathrm{RWA}}=P_\B\left[H_0+D\sum_m
       (u_ma_m\sigma_++u_m^*a_m^\dagger\sigma_-)\right]P_\B.
 \label{eq:rwa}
\end{equation}
With $\mathcal N_f=\sum_m n_m$, $\mathcal N=\mathcal N_f+\ket e\bra e$
and $\Pi=(-1)^{\mathcal N}$,
\begin{equation}
 [H_{\B,\mathrm{RWA}},\mathcal N]=0,\qquad [H_\B,\Pi]=0.
 \label{eq:symmetries}
\end{equation}
RWA transitions change photon and atomic excitation by opposite units.
Counter-rotating terms change their sum by $\pm2$, preserving its parity.
Occupation projectors commute with the diagonal symmetry operators, so these
identities survive all three projections.

For the retained zero mode and \code{sqrt},
\begin{equation}
 \omega(0)=f(0)=0,\qquad [H_\B,n_0]=0.
 \label{eq:zero}
\end{equation}
It can carry initial norm without free energy or direct TLS coupling.
In full, fixing $n_0=r$ leaves the same nonzero-mode Hamiltonian for each $r$.
In total-cap the corresponding sector instead obeys
\begin{equation}
 \sum_{m:k_m\ne0}n_m\leq N-r.
\end{equation}
A conserved zero-mode occupation thus changes the available nonzero photon
space. In truncated, positive zero-mode occupation cannot coexist with
occupation of another mode. Deleting zero changes preparation normalization
and these finite spaces even though $f(0)=0$. For flat, $f(0)=1$ and the zero
mode couples, so $n_0$ is generally no longer conserved. At $N=0$ all interactions
are projected away. At explicit $M=1$ with sqrt the only mode is zero and
only TLS free evolution remains.

\paragraph{Propagation, output sampling and numerical normalization.}
The engine solves the time-independent equation
\begin{equation}
 \ii\frac{d}{dt}\ket{\psi(t)}=H_\B\ket{\psi(t)},\qquad
 \ket{\psi(t)}=e^{-\ii H_\B t}\ket{\psi_0}.
 \label{eq:evolution}
\end{equation}
The default integration method used by this engine is BDF. Relative and
absolute tolerances are \code{rtol=1e-9} and \code{atol=1e-11}.
Output renormalization is disabled (\code{normalize_output=False}). Output \code{dt}
is not the adaptive solver's internal step. Nominal outputs are
$0,\code{dt},\ldots,J\code{dt}$, $J=\lfloor T/\code{dt}\rfloor$,
with $T$ appended if needed. Endpoint roundoff is handled at $10^{-12}$;
both zero and exact $T$ are always included. If \code{dt} exceeds $T$, the
list is $[0,T]$.

Exact evolution conserves $\|v\|^2$ and $\langle H_\B\rangle$. The reported
norm is the squared norm, not its square root. Coefficients, populations and
energies are not renormalized after integration. Entropy and fidelity normalize
internally. If $v_{\mathrm{num}}=(1+\epsilon)e^{\ii\theta}v_{\mathrm{exact}}$,
its fidelity is one after normalization but its raw norm is $|1+\epsilon|^2$.
Fidelity alone therefore does not certify norm accuracy.

With \code{store_state=True}, solver history and coefficient arrays contain
all outputs. With \code{store_state=False}, only the final ket and final
coefficient arrays are retained; the output-time list still exists.
Time-series fidelity requires both histories. Diagnostic methods accept
\code{t=None} for all retained states, \code{t=-1} for final, or a physical
time in $[0,T]$. With history, a time selects the nearest output without
interpolation; a final-only run rejects earlier times. The wavefunction
method instead takes an output \emph{index}.

\subsection{Observables, exact fidelities and physical interpretation}
\paragraph{TLS reduction and entanglement.}
Tracing equation~\eqref{eq:index} over field configurations gives
\begin{equation}
 \rho_a=C^{\mathsf T}C^*,\qquad
 (\rho_a)_{ss'}=\sum_i C_{i,s}C_{i,s'}^*,\qquad
 P_e=\sum_i|C_{i,e}|^2,\qquad \Tr\rho_a=\|v\|^2.
 \label{eq:rho}
\end{equation}
The transpose is intentionally not an adjoint: rows of $C$ are photon
configurations and columns are TLS labels. For
$\widehat\rho_a=\rho_a/\|v\|^2=\begin{psmallmatrix}p_g&z\\z^*&p_e\end{psmallmatrix}$,
\begin{equation}
 \lambda_\pm=\frac{1\pm\sqrt{(p_g-p_e)^2+4|z|^2}}2,\qquad
 S_a=-\sum_{r=\pm}\lambda_r\log\lambda_r.
 \label{eq:entropy}
\end{equation}
The logarithm is natural and $0\log0=0$, so $0\leq S_a\leq\log2$.
For a global pure ket this measures field--TLS entanglement: zero for a
product, maximal for a maximally mixed TLS reduction. The excited-probability
method returns the raw diagonal of $\rho_a$, exposing normalization drift.

\paragraph{Photon sectors and Fourier reconstruction.}
Directional populations refer specifically to one photon and ground-state TLS:
\begin{align}
 P_{+,1g}&=\sum_{m:k_m>0}|C_{i(\bm e_m),g}|^2,&
 P_{-,1g}&=\sum_{m:k_m<0}|C_{i(\bm e_m),g}|^2,\\
 P_{0,1g}&=\sum_{m:k_m=0}|C_{i(\bm e_m),g}|^2,&
 p_\nu&=\sum_{i:\sum_m n_{i,m}=\nu}\sum_s|C_{i,s}|^2.
 \label{eq:sectors}
\end{align}
Their keys are \code{p_transmitted}, \code{p_reflected}, \code{p_zero_1g}
and \code{p_<nu>_photon}. Exact sum rules are
\begin{equation}
 P_{+,1g}+P_{-,1g}+P_{0,1g}=P_{1g}\leq p_1,\qquad
 \sum_\nu p_\nu=\|v\|^2.
 \label{eq:populationchecks}
\end{equation}
The difference $p_1-P_{1g}$ is the one-photon, excited-TLS population.
Photon totals range through $MN$ in full and $N$ in the other two bases.
Directional populations vanish at $N=0$.

The extracted function is
\begin{equation}
 \phi_{1g}(x,t)=\frac1{\sqrt L}\sum_m C_{i(\bm e_m),g}(t)e^{\ii k_mx},\qquad
 \int_{x_*}^{x_*+L}|\phi_{1g}(x,t)|^2dx=P_{1g}(t).
 \label{eq:wavefunction}
\end{equation}
This is a component of the ket, not a density of all photons in a multiphoton
state. Its integrated weight may be less than one. The integral uses
orthogonality from equation~\eqref{eq:grid}; a plotting mesh approximates it.
The function is periodic. Under free evolution a positive-momentum component
obeys $\phi(x,t)=\phi(x-t,0)$ and a negative-momentum component obeys
$\phi(x,t)=\phi(x+t,0)$. Both directions can be present in the packet.

The imposed spatial phases give, for an input in the $1g$ sector,
\begin{align}
 (H\psi)_{\bm0,e}
 &=\frac{\ii D}{\sqrt L}\sum_m
           f_mC_{i(\bm e_m),g}e^{-\ii k_mx_{\mathrm{tls}}},\\
 \text{for \code{flat}:}\quad
 (H\psi)_{\bm0,e}&=\ii D\phi_{1g}(-x_{\mathrm{tls}}).
 \label{eq:position}
\end{align}
For sqrt the same coordinate enters the field weighted by $\sqrt{|k_m|}$.
The scattering plot therefore marks $-x_{\mathrm{tls}}$ in its positive-Fourier
coordinate. The parameter remains the one in the requested annihilation
phase. Identifying it directly with the plotted coordinate would require
the opposite annihilation-phase sign; at zero they coincide. A general
non-RWA input can have additional contributions to this target from other sectors.

\paragraph{Total Hamiltonian expectation and energy conservation.}
For the joint coefficient vector $v_q(t)=C_{i,s}(t)$, $q=2i+s$, and matrix
$H_{rc}=\bra rH_\B\ket c$, the total energy of an exactly normalized state is
\begin{equation}
 E(t)=\bra{\psi(t)}H_\B\ket{\psi(t)}
     =v(t)^\dagger H_{\B}v(t)
     =\sum_{r,c}v_r^*(t)H_{rc}v_c(t).
 \label{eq:total-energy}
\end{equation}
Here $r,c$ label joint field--TLS basis states. Since $H_\B$ is Hermitian,
the sum is real. With $\dot\psi=-\ii H_\B\psi$ and $\partial_tH_\B=0$,
\begin{equation}
 \frac{dE}{dt}
 =\bra{\dot\psi}H_\B\ket\psi+\bra\psi H_\B\ket{\dot\psi}
 =\ii\bra\psi H_\B^2\ket\psi-\ii\bra\psi H_\B^2\ket\psi=0.
 \label{eq:total-energy-conservation}
\end{equation}
Thus $E(t)=E(0)$ for exact evolution in each finite model. Energy conservation
is a numerical diagnostic, not a certificate of full-state convergence.

The method \code{compute_energy} returns the raw quadratic expression, without
renormalizing propagated kets. If the numerical squared norm is
$\eta(t)=v(t)^\dagger v(t)$, distinguish
\begin{equation}
 E_{\mathrm{raw}}(t)=\Re[v(t)^\dagger H_{\B}v(t)],\qquad
 E_{\mathrm{normalized}}(t)=E_{\mathrm{raw}}(t)/\eta(t)\quad(\eta>0).
\end{equation}
They coincide when $\eta=1$. For an already propagated experiment:
\begin{lstlisting}
observables = experiment.compute_observables()
time = observables['time']
E_raw = experiment.compute_energy()
E_normalized = E_raw / observables['norm']
\end{lstlisting}
The returned energy array spans the retained outputs. An explicit time selects
the nearest stored output; a final-only run returns its final value. Monitoring
$\max_j|E_{\mathrm{raw}}(t_j)-E_{\mathrm{raw}}(0)|$ alongside the norm and solver
tolerances tests conservation without dividing by a possibly zero initial energy.

\paragraph{Implicit embedding in a common full space.}
Comparisons require a common physical box and consistent conventions.
Let $\K_*=\K_A\cup\K_B$ and $N_*\geq\max(N_A,N_B)$; the ambient space is
\begin{equation}
 \mathcal H_* = \operatorname{span}
 \{\ket{\bm n;s}:0\leq n_m\leq N_*,\ k_m\in\K_*,\ s=g,e\}.
 \label{eq:ambient}
\end{equation}
A physical mode map $\mu_X$ defines
\begin{equation}
 (\iota_X\bm n)_{\mu_X(m)}=n_m,\qquad
 (\iota_X\bm n)_\ell=0\text{ for absent modes},\qquad
 \iota_X\ket{\bm n;s}=\ket{\iota_X\bm n;s}.
\end{equation}
Distinct tuples stay orthogonal, so
\begin{equation}
 \iota_X^\dagger\iota_X=I_X,\qquad
 \|\iota_X\psi_X\|=\|\psi_X\|.
 \label{eq:isometry}
\end{equation}
For example, adding $k=0$ maps $\ket{n_{-1},n_{+1};s}$ to
$\ket{n_{-1},0,n_{+1};s}$. Positive occupation on that added mode is
orthogonal to every such embedded ket. Reordering existing modes only permutes
coordinates and cannot change an overlap.

For individually normalized kets,
\begin{align}
 A_{AB}&=\braket{\iota_A\psi_A}{\iota_B\psi_B}
       =\sum_{\bm n\in\iota_A\B_A\cap\iota_B\B_B}
                       \sum_s C_{A,\bm n,s}^*C_{B,\bm n,s},\\
 F_{\mathrm{state}}&=|A_{AB}|^2.
 \label{eq:fstate}
\end{align}
The intersection enters because missing amplitudes are zero. No projection
and renormalization on the intersection is performed. For orthogonal
configurations $x,y$, take $\ket{\psi_A}=\ket x$ and
$\ket{\psi_B}=\sqrt p\ket x+\sqrt{1-p}\ket y$. Then
\begin{equation}
 F_{\mathrm{state}}=p;
 \qquad\text{renormalizing the common component would incorrectly give }1.
 \label{eq:intersectionexample}
\end{equation}
For distinct one-photon modes $a,b$, the states
$(\ket{0;g}+\ket{1_a;g})/\sqrt2$ and
$(\ket{0;g}+\ket{1_b;g})/\sqrt2$ share only vacuum. Their overlap is $1/2$
and their state fidelity is $1/4$.

Keys contain sorted pairs $(\text{physical mode index},n_m)$ for $n_m>0$;
vacuum has the empty key. Each stores two TLS amplitudes. All shared
configurations contribute, including several occupied modes. No vector of
dimension $2(N_*+1)^{|\K_*|}$ is created. The union retains the first grid's
order and appends absent modes; sorting keys makes that internal order irrelevant.
Momentum matching uses absolute tolerance $10^{-12}$, zero relative tolerance,
and rejects ambiguous identifications.

The time-series interface checks equal $L$ to absolute tolerance $10^{-12}$
and exactly equal output arrays before pairing states. Low-level
\code{compare_states} has no $L$ or Hamiltonian input, so its caller must
ensure physical comparability. Matching numerical momenta from different
boxes alone does not certify identical normalized modes.

\paragraph{Atomic fidelity, normalization and bounds.}
The atomic metric is squared Uhlmann fidelity,
\begin{equation}
 F_{\mathrm{atom}}=
 \left[\Tr\sqrt{\sqrt{\widehat\rho_{a,A}}\,
                    \widehat\rho_{a,B}\sqrt{\widehat\rho_{a,A}}}\right]^2.
 \label{eq:fatom}
\end{equation}
The code squares the value returned by \code{qutip.fidelity}. For qubits,
\begin{equation}
 F_{\mathrm{atom}}=\Tr(\widehat\rho_{a,A}\widehat\rho_{a,B})
                  +2\sqrt{\det\widehat\rho_{a,A}\det\widehat\rho_{a,B}}.
 \label{eq:qubit}
\end{equation}
To derive this, let $X=\sqrt\rho\,\sigma\sqrt\rho$ have eigenvalues $x_1,x_2$.
Then $(\Tr\sqrt X)^2=x_1+x_2+2\sqrt{x_1x_2}$,
$\Tr X=\Tr(\rho\sigma)$ and $\det X=\det\rho\det\sigma$.
Rank-deficient cases follow by continuity. For Bloch vectors $\bm r,\bm s$,
\begin{equation}
 F_{\mathrm{atom}}=\frac{1+\bm r\cdot\bm s
              +\sqrt{(1-|\bm r|^2)(1-|\bm s|^2)}}2.
\end{equation}
It reduces to squared overlap for pure atomic states, and against a pure
$\ket e$ to the normalized excited population. Adding vacuum modes leaves
the atomic reduction invariant; no ambient allocation is required.

Both metrics normalize each entire finite nonzero ket, not a common subset.
Values are clipped to $[0,1]$ for roundoff overshoots. Raw norm diagnostics
remain visible. Monotonicity under the common partial trace yields
\begin{equation}
 0\leq F_{\mathrm{state}}\leq F_{\mathrm{atom}}\leq1.
 \label{eq:monotonicity}
\end{equation}
For instance, $\ket{1_+;g}$ and $\ket{1_-;g}$ have state fidelity zero
and atomic fidelity one. No ratio is calculated: a quotient of these metrics
is not generally a fidelity of the field reductions.

For normalized embedded pure states, $\rho_A-\rho_B$ has nonzero eigenvalues
$\pm\sqrt{1-|\braket{\psi_A}{\psi_B}|^2}$ in their two-dimensional span.
Therefore, for any common projector, or more generally $0\leq Q\leq I$,
\begin{equation}
 |\Tr[Q(\rho_A-\rho_B)]|
 \leq\tfrac12\|\rho_A-\rho_B\|_1=\sqrt{1-F_{\mathrm{state}}}.
 \label{eq:bound}
\end{equation}
This bounds disagreement between normalized populations of two finite
calculations, not accuracy against a continuum experiment. For TLS-only
measurements an analogous upper bound uses $\sqrt{1-F_{\mathrm{atom}}}$.

\paragraph{Sample means and scattering interpretation.}
The coupling sweep and mode selection report
\begin{equation}
 \overline F_{\mathrm{samples}}=\frac1{N_t}\sum_{j=0}^{N_t-1}F(t_j).
 \label{eq:mean}
\end{equation}
This differs from final fidelity and a time integral. A trapezoidal time
average would instead be
\begin{equation}
 \overline F_{\mathrm{time}}\simeq\frac1T\sum_{j=0}^{N_t-2}
       \frac{F(t_j)+F(t_{j+1})}{2}(t_{j+1}-t_j).
 \label{eq:quadrature}
\end{equation}
Only equation~\eqref{eq:mean} is implemented. Endpoints have different weights
even on a uniform grid; a shorter final interval still has an equally weighted
output in the sample mean. Initial metrics are stored separately. Cap convergence
compares final kets for adjacent entries in the supplied increasing cap list,
which need not differ by one.

Scattering supplies finite-packet sector populations and Fourier profiles.
If $\sum_{k_m<0}|c_m|^2>0$, final negative-momentum population is not by itself
a reflection probability for an exclusively positive-momentum input.
Subtracting the initial population need not remove interference.
For RWA evolution from one photon and ground TLS, excitation conservation gives
\begin{equation}
 P_{+,1g}+P_{-,1g}+P_{0,1g}+P_e=1
 \label{eq:rwasectors}
\end{equation}
for a normalized state. Without RWA additional sectors can carry norm, so
this four-term identity is not general.

Extracting $t(k)$ or $r(k)$ needs identifiable outgoing packets, control of
periodic returns and a phase convention relative to free propagation.
Counter-rotating terms can require dressed asymptotic states and additional
photon channels. The calculated quantities remain finite-time projections;
neither fidelity is an $S$-matrix element.

\section{Notebook control, resources and source-line commentary}
\subsection{Shared configuration and callable Python experiments}
\paragraph{Configuration and execution.}
\code{ExperimentConfig} collects three physical dictionaries, cutoffs, cap,
basis name, RWA, storage choice, grid control and selection windows. Validation
occurs when resolving the grid, building the basis/Hamiltonian and preparing
the ket. \code{with_D} returns a copy with a new atomic dictionary.
An \code{Experiment} copies its physical dictionaries before construction,
so sweeps do not overwrite the supplied physical configuration.
A notebook parameter cell can use:
\begin{lstlisting}
import numpy as np
from dataclasses import replace
from src.xp_config import ExperimentConfig
from src.experiment import estimate_config
from src.fidelities import fidelities_over_time
from experiment.atom_evolution import (
    run_atom_evolution, plot_atom_evolution,
)

config = ExperimentConfig(
    param_photon={'k_0': 1., 'sigma_k': .4, 'x_0': -2.,
                  'state': 'number', 'n': 1},
    param_atom={'omega_0': 1., 'D': .2, 'L': 2*np.pi,
                'x_tls': 0., 'coupling': 'sqrt',
                'initial_state': 'g'},
    param_time_evol={'T': 4., 'dt': .05,
                     'rtol': 1e-9, 'atol': 1e-11},
    n_max=2, CTRL_M_EXPLICIT=True, M=5,
    truncation='full+totalcap', store_state=True)

# Estimate every basis that will actually be propagated.
for scheme in ('full+totalcap', 'truncated'):
    estimate_config(replace(config, truncation=scheme))
runs = run_atom_evolution(config)
metrics = fidelities_over_time(
    runs['truncated'], runs['full+totalcap'])
fig = plot_atom_evolution(runs)
\end{lstlisting}
For cutoff control, use \code{CTRL_M_EXPLICIT=False} and
\code{cutoffs={'ir_cutoff': 0., 'uv_cutoff': 2.}} instead.
At this $L$ the base grid is identical. Direct execution constructs
\code{Experiment}, calls \code{propagate_state}, then \code{compute_observables};
the common experiment helper performs those three operations.

\paragraph{Five experiments and their output contracts.}
All four migrated Uri notebooks explicitly set \code{coupling='sqrt'}.
The scattering notebook remains alongside them. Modules do not simulate
or change the Matplotlib backend on import. Numerical execution and plotting
are separate functions, with plotting returning figure objects.
\begin{longtable}{p{.34\textwidth}p{.57\textwidth}}
\toprule Module & Calculation and main return value\\\midrule\endhead
\code{atom_evolution.py} & Dictionary of propagated \code{Experiment}
objects by requested basis. Forces history storage; plots $P_e$ and entropy.
The notebook additionally calculates the two fidelities.\\
\code{full_cumulative_fidelity.py} & Dictionary of increasing caps,
retained final-only runs, and initial/final metrics for adjacent cap-list
entries. Always uses total-cap.\\
\code{sweep_D_fidelity.py} & Dictionary of physical couplings, candidate
names, comparator name, sample means and initial metrics. Forces history.\\
\code{compare_mode_selection.py} & Coupling sweeps with/without selection,
and window heatmaps at fixed coupling. Requires history. Comparator is
always unselected total-cap.\\
\code{scattering.py} & One experiment with the user's dictionary API;
optional observable CSV. Plots $1g$ profiles and directional/zero/TLS/norm
populations.\\
\bottomrule
\end{longtable}
With $S$ candidate schemes, $Q$ coupling values, $C$ caps, and $P,A$ window
factor counts, the layouts are:
\begin{longtable}{p{.28\textwidth}p{.63\textwidth}}
\toprule Result & Array shapes and axes\\\midrule\endhead
Atom evolution & Ground/excited coefficient arrays $(N_t,B)$; population,
entropy and time-series fidelity arrays $(N_t,)$.\\
Cap convergence & Each metric and its \code{_initial} counterpart:
$(C-1,)$. Entry $j$ compares supplied caps $j$ and $j+1$.\\
Coupling sweep & Each metric and its \code{_initial} counterpart:
$(S,Q)$, scheme before coupling. Unsuffixed metrics are sample means.\\
Selection & Sweep metrics $(S,2,Q)$, with middle axis
$[\mathrm{False},\mathrm{True}]$. Each \code{_heatmap}: $(S,P,A)$,
photon window before atomic window.\\
Scattering & Observable arrays of length $N_t$ with history, one otherwise;
wavefunction array indexed by the supplied spatial samples.\\
\bottomrule
\end{longtable}
For final-only runs the coefficient arrays have shape $(B,)$, not $(1,B)$;
observable methods promote them internally as needed. Diagnostics with
\code{t=None} return a length-one array/list then, whereas an explicit final
time returns one scalar or matrix.

Cap convergence compares equal physical/grid inputs with different caps.
Its default CLI preparation is coherent, so initial projection differences
are expected. The coupling sweep defaults to a propagated total-cap
comparator; \code{reference} can select another finite basis. This is distinct
from the ambient full space in equation~\eqref{eq:ambient}. A candidate with
the comparator's scheme reuses its trajectory.

The selection study rebuilds an unselected comparator at each $D$. Heatmaps
use \code{D_heatmap} when given, otherwise the input configuration's coupling.
Their returns contain only the sample means, not separate initial heatmap
metrics; the sweep arrays do contain the \code{_initial} values. This distinction
matters when separating selected-preparation and dynamical effects.

\paragraph{Command-line execution, exports and notebook state.}
The five experiments are invoked through their modules in \code{experiment/}.
Shared arguments use \code{--option value}; old numerical values still require
the convention conversions given in the comparative analysis. From the
repository root:
\begin{lstlisting}[language={}]
python -m experiment.scattering --D 0.2 --ctrl-m-explicit --M 5
python -m experiment.sweep_D_fidelity --D-values 0.02,0.1,0.2 \
    --out results/sweep.png
python -m experiment.atom_evolution --D 0.2 --out results/atom.png
\end{lstlisting}
The CLI defaults to cutoff control; \code{--M} is active only with
\code{--ctrl-m-explicit}. Both profiles are available, with default sqrt;
\code{--RWA} enables RWA. Solver options can be set in the Python time
dictionary. Notebooks resolve the root whether the kernel starts there or in
\code{notebooks/}. Markdown, comments, labels and function docstrings are in
English. Each function defines the meaning, types and array shapes of its
inputs and outputs, with formulas for the relevant numerical operations.

Resource and diagnostic outputs are stored in pandas DataFrames rather than
line-by-line prints. Numerical estimate keys remain available alongside
\code{grid_table} and \code{resource_table}; setting \code{print_report=False}
suppresses automatic display without discarding either table. The notebooks
keep named DataFrames for reuse or export:
\begin{lstlisting}
estimate = estimate_config(config, print_report=False)
grid_table = estimate['grid_table']
resource_table = estimate['resource_table']
display(grid_table, resource_table)

# After propagation:
observable_history = experiment.observables_dataframe()
observable_summary = experiment.observables_dataframe(summary=True)
display(observable_summary)
\end{lstlisting}
The history has physical time as its index and raw observables as columns.
The summary has observable names as rows and \code{initial}/\code{final}
columns, including a row for time. Final-only storage produces just a
\code{final} column: its retained state is never presented as an initial
sample. Neither conversion normalizes populations nor rounds stored values.
Fidelity DataFrames expose initial, final or sample-mean metrics with explicit
labels; sweeps use indices naming the scheme, coupling, selection flag or
compared caps. Window scans use scheme and both window radii as index levels.
DataFrames support their usual CSV export. The estimator records excluded
memory costs in \code{resource_table.attrs['memory_excludes']} and the
heuristic feasibility rule in \code{resource_table.attrs['feasibility_rule']}.

With \code{--out}, a CLI saves a figure and same-stem NPZ with the supplied
arrays. Its metadata field is \code{configuration_json}.
JSON stores the base configuration;
additional sweep choices and derived grids must be read from accompanying
arrays or reconstructed. Complex values are represented as strings.
Ordinary exported arrays require no pickle:
\begin{lstlisting}
import json
with np.load('results/sweep.npz', allow_pickle=False) as data:
    metadata = json.loads(str(data['configuration_json']))
    D_values = data['D']
    F_state_mean = data['F_state']
\end{lstlisting}
The atom-evolution CLI exports time and excited populations; entropy and
notebook fidelity curves are not automatically included in that export.
Cap exports omit Python run objects. Optional scattering CSV writes
\code{results/scattering.csv}, overwriting a prior requested export, with
one column per observable and no configuration metadata.

\subsection{Correspondence between controls and resource costs}
\paragraph{What the grid report establishes.}
The same \code{resolve_grid} serves propagation and estimation before Fock
allocation. Base and selected descriptions contain counts, signed extrema,
radial extrema, zero presence, smallest nonzero momentum and integer indices.
For integer index set $I$, the report tests
\begin{equation}
 I=\{j:-j_{\max}\leq j\leq j_{\max},\ |j|\geq j_{\min}\},\quad
 j_{\min}=\min_{j\in I}|j|,\quad j_{\max}=\max_{j\in I}|j|.
 \label{eq:radialexact}
\end{equation}
This is \code{radial_cutoffs_exact}. If also $j_{\min}=0$, the count is
\code{equivalent_M}; otherwise that field is \code{None}. Extrema alone
cannot represent asymmetry or holes. At $L=2\pi$:
\begin{center}
\begin{tabular}{l l c c}
\toprule Inputs & Grid & Radial exact & Equivalent $M$\\\midrule
$M=5$ & $\{-2,-1,0,1,2\}$ & Yes & 5\\
IR/UV $=(0,2)$ & $\{-2,-1,0,1,2\}$ & Yes & 5\\
IR/UV $=(0.2,2.8)$ & $\{-2,-1,1,2\}$ & Yes & None\\
IR/UV $=(0,0)$ & $\{0\}$ & Yes & 1\\
Selected $\{-2,0,2\}$ & $\{-2,0,2\}$ & No & None\\
Selected $\{0,1,2\}$ & $\{0,1,2\}$ & No & None\\\bottomrule
\end{tabular}
\end{center}
The first two grids give identical finite dynamics at equal other inputs.
A positive-IR band has no allowed equivalent explicit count. The report
states this rather than proposing a count that silently inserts zero.

\paragraph{Vector storage and exact dimension growth.}
A complex128 ket occupies $16d$ bytes. With $S=\code{store_state}$,
\begin{align}
 G_{\mathrm{ket}}&=\frac{16d}{2^{30}},&
 G_{\mathrm{history}}&=G_{\mathrm{ket}}
        \begin{cases}N_t,&S,\\1,&\neg S,\end{cases}\\
 G_{\mathrm{vectors}}&\simeq2G_{\mathrm{history}},&
 \operatorname{nnz}(H)&\leq d(1+2M).
 \label{eq:memory}
\end{align}
GiB means $2^{30}$ bytes. The factor two estimates solver kets and the
contiguous coefficient array. Alternating ground/excited views share that
array and do not add two more full copies. The initial ket, QuTiP objects,
transient conversions and solver workspace can increase storage further.

For $M=5,N=2$, joint dimensions are $22,42,486$ for truncated, total-cap and
full. At $T=4,\code{dt}=0.05$, $N_t=81$ and estimated retained vector bytes
are $57\,024$, $108\,864$, and $1\,259\,712$. At fixed $N$ and large $M$,
\begin{equation}
 d_{\mathrm{total}}\sim2M^N/N!,\qquad
 d_{\mathrm{truncated}}\sim2NM,\qquad d_{\mathrm{full}}=2(N+1)^M.
\end{equation}
Polynomial growth at fixed cap does not make high caps cheap; the exact
binomial dimension should be evaluated before allocation.

\paragraph{Sparse transitions and excluded costs.}
Let $A_\B$ count photon edges $\bm n\leftrightarrow\bm n+\bm e_m$ in a
basis, without TLS labels. For $N\geq1$,
\begin{align}
 A_{\mathrm{truncated}}&=MN,\\
 A_{\mathrm{total}}&=M\binom{M+N-1}{N-1},\\
 A_{\mathrm{full}}&=MN(N+1)^{M-1}.
 \label{eq:edges}
\end{align}
For total-cap, remove one photon in the changed mode and count configurations
with total at most $N-1$. For full, choose the mode, one of its $N$ adjacent
levels, and the other $M-1$ occupations. At $N=0$ all counts are zero.
The engine assembles $4A_\B$ directed TLS-flipping triplets without RWA,
$2A_\B$ with it. Zero coefficients, such as the sqrt zero mode, may still
occupy recorded triplet slots. Physically nonzero matrix entries can be
fewer, especially at $D=0$. The estimator uses the simpler bound above.

A complex CSR matrix with $Z$ stored entries uses approximately
\begin{equation}
 16Z+b_{\mathrm{idx}}Z+b_{\mathrm{idx}}(d+1)\quad\text{bytes}
 \label{eq:csrbytes}
\end{equation}
for its numerical arrays, where $b_{\mathrm{idx}}$ depends on the sparse
index type. The engine retains $H_0$, $V$ and its QuTiP Hamiltonian.
Temporary row/column/value lists and COO/CSR conversion add assembly peaks. Basis tuples, dictionary entries and
photon-number arrays have at least $MB$ occupation entries plus object
overhead. These costs are excluded from the returned vector estimate.

\code{feasible} is only the heuristic $d\leq100000$ and
$G_{\mathrm{vectors}}\leq1$ GiB. It is not a memory guarantee, a physical
accuracy criterion or an automatic engine limit. Notebooks check it before
allocation. Estimates concern one trajectory: retained objects at several
caps/schemes add their costs. Cap convergence keeps all its runs, even though
they are final-only. Sweep peaks include a comparator and candidate together.
The largest individual estimate is therefore not a total collection estimate.

Implicit fidelity keys examine $O(M_AB_A+M_BB_B)$ occupation entries.
Sorting each key adds work dependent on its occupied-mode count; matching
physical modes searches the current union for each second-grid mode.
Keys and dictionaries add memory, but no ambient full vector is allocated.
Choosing full for an actual propagation still incurs its exponential dimension.

\subsection{Numbered code and comments on individual operations}
Listings below are loaded directly from executed sources. Commentary tables
identify current line ranges and relate their operations to the equations.
The generator finds ordered engine markers and syntax-tree function ranges
for experiments; line numbers are not copied from old versions. Comments
cover validations, tensor ordering, phases, missing transitions, normalization,
output shapes and statistical choices. Core-engine ranges are deliberately
narrower than experiment function ranges.
To reproduce the documentation from the repository root:
\begin{lstlisting}[language={}]
python documentation/build_code_notes.py
latexmk -pdf -cd documentation/main.tex
\end{lstlisting}
These commands generate documentation only.
% Generated by build_code_notes.py; edit the generator.
\paragraph{\code{src/xp_config.py}}
\lstinputlisting{../src/xp_config.py}
\begin{longtable}{p{.12\textwidth}p{.79\textwidth}}
\toprule Lines & Mathematical and implementation commentary\\\midrule\endhead
1--5 & This module defines configuration data only. Imports do not allocate a grid or launch propagation.\\
6--7 & The dataclass preserves named inputs and supports \code{dataclasses.replace}; validation occurs in the consuming numerical routines.\\
8--8 & Packet inputs $k_0,\sigma_k,x_0$, preparation type, and either integer $n$ or complex $\alpha$; see \eqref{eq:packet}--\eqref{eq:coherent}.\\
9--9 & $\omega_0,D,L,x_{\mathrm{tls}}$, coupling profile, and optional initial TLS label. The box length is stored here but also controls the photon grid.\\
10--10 & Final time $T$, output spacing \code{dt}, and optional integration method/tolerances. Output spacing does not fix internal solver steps.\\
11--11 & The radial band is required only under cutoff control. Explicit-count control ignores this dictionary.\\
12--15 & $N$ and the basis select the retained occupation set, \eqref{eq:bases}. RWA selects Hamiltonian terms; storage selects retained outputs, not a different Hilbert space.\\
16--17 & True selects exact odd $M$ with zero, \eqref{eq:explicit}; false selects the radial band, \eqref{eq:cutoffs}. No hidden decrement of $M$ occurs.\\
18--21 & Additional mode windows, applied after resolving the base grid. Their radii are factors times $\sigma_k$, \eqref{eq:selection}.\\
22--24 & Return a replaced configuration with a new atomic dictionary. Other inputs are reused; the supplied dictionary is not overwritten during a sweep.\\
\bottomrule\end{longtable}
\paragraph{\code{src/grid.py}}
\lstinputlisting{../src/grid.py}
\begin{longtable}{p{.12\textwidth}p{.79\textwidth}}
\toprule Lines & Mathematical and implementation commentary\\\midrule\endhead
1--13 & \code{integer} accepts integral types above its minimum, including NumPy integers. It rejects booleans and all floats rather than silently coercing them.\\
14--15 & One grid generator for both evolution and estimation; its inputs describe the base grid, before any optional windows.\\
16--20 & Convert and require finite $L>0$. Require an actual boolean control so a truthy string cannot silently select the wrong branch.\\
21--21 & Compute the physical spacing $\Delta k=2\pi/L$, \eqref{eq:grid}; it is never inferred from a chosen number of modes.\\
22--28 & Require positive odd integer $M$ and construct indices $-J,\ldots,J$, $J=(M-1)/2$. Even counts raise an error before basis allocation.\\
29--33 & Cutoff control requires finite ordered radial bounds $0\leq\mathrm{IR}\leq\mathrm{UV}$. These validation steps are skipped when $M$ controls the grid.\\
34--35 & Ceiling/floor of cutoff ratios with dimensionless $10^{-12}$ boundary tolerances. The positive lower index is subsequently restricted to at least one; compare \eqref{eq:count}.\\
36--38 & Enumerate positive indices, reflect them in reversed order for the negative side, and add zero only when input IR is exactly zero. A tiny positive IR is not rounded to zero.\\
39--44 & Multiply integer indices by $\Delta k$ and reject an empty band. Sorted signed momenta, not frequencies, are returned.\\
45--51 & Validate positive finite $\sigma_k$ and nonnegative finite factors. This is a requested subset operation; it is not a right-moving preparation filter.\\
52--54 & Window centres $k_0,+\omega_0,-\omega_0$ with radii $w_p\sigma_k,w_a\sigma_k$, as in \eqref{eq:selection}.\\
55--63 & For each window use an absolute momentum tolerance, add the nearest mode if empty, and union its mask. Ties use the first sorted entry. The returned mask preserves base-grid ordering.\\
64--72 & Resolve the same base grid for estimator and engine. With selection return its indexed subset; otherwise return a distinct array copy with the same values.\\
73--77 & Recover integer indices from the physical spacing and compute radial extrema. This representation report is independent of the photon cap and basis type.\\
78--80 & Test equality with the entire symmetric radial band, \eqref{eq:radialexact}; minima/maxima alone would miss holes or unpaired endpoints. Compute the smallest nonzero $|k|$ separately.\\
81--88 & Report count, signed extrema, effective cutoffs, zero, integer indices and exact representability. An equivalent explicit $M$ exists only for a radial-exact set containing zero.\\
89--89 & Return separate base and selected descriptions with their common physical spacing. A selected count may differ from the explicit base-grid input.\\
\bottomrule\end{longtable}
\paragraph{\code{src/states.py}}
\lstinputlisting{../src/states.py}
\begin{longtable}{p{.12\textwidth}p{.79\textwidth}}
\toprule Lines & Mathematical and implementation commentary\\\midrule\endhead
1--11 & Combinations with replacement enumerate total-cap occupations; a Cartesian product enumerates full. These algorithms avoid testing every candidate product against a total cap.\\
12--16 & Validate at least one retained mode and nonnegative integer cap. Store the representation name; $N=0$ is a valid vacuum-only photon space.\\
17--17 & One zero tuple of length $M$ represents the photon vacuum. TLS components are attached later through the alternating index rule \eqref{eq:index}.\\
18--24 & Append $n\bm e_m$ in mode-first, occupation-second order after vacuum. Multimode configurations are absent at every cap; dimension $2(1+MN)$.\\
25--29 & For each total from zero to $N$, enumerate multisets of occupied mode indices. \code{bincount} produces the tuple without duplicates; dimension $2\binom{M+N}{N}$.\\
30--33 & Cartesian product of $0,\ldots,N$ on each mode. The final occupation changes fastest. Unknown names, including the removed extra-oscillator scheme, raise an error.\\
34--39 & Store the exact tuples, their dictionary indices and photon totals. Ket dimension is twice the configuration count because the TLS remains a separate binary factor.\\
40--46 & Read finite signed $k_0$, finite $x_0$ and positive finite $\sigma_k$. All values in the supplied selected momentum array enter preparation.\\
47--49 & Gaussian amplitude and phase from \eqref{eq:packet}. Subtract the largest real exponent, then normalize the whole packet; the common stabilizing factor cancels.\\
50--55 & Define the normalized $g,e,+,-$ vectors in TLS order $(g,e)$ and reject unknown initial labels. Ground is the default.\\
56--58 & Allocate the joint complex ket and select number or coherent preparation. The default is number, with $n=1$ below.\\
59--66 & Require integer $0\leq n\leq N$. For $n=0$ write the TLS vector onto vacuum exactly once, independent of how many modes exist.\\
67--71 & For $n>0$, place $c_m\ket{a_0}$ at occupation $n\bm e_m$ in alternating components. This is \eqref{eq:number}, not the collective-mode Fock state \eqref{eq:collective}.\\
72--75 & Read and validate complex $\alpha$. The algorithm projects a product of coherent states, not a single coherent state restricted to one ray.\\
76--82 & Multiply $(\alpha c_m)^{n_m}/\sqrt{n_m!}$ over all modes for each allowed tuple, then multiply both TLS components. The common coherent exponential is omitted because it cancels in \eqref{eq:coherent}.\\
83--86 & Normalize the complete projected ket once, rejecting zero or nonfinite norm. The retained coherent mass and its basis dependence are derived in \eqref{eq:weights}; no intersection normalization is involved.\\
\bottomrule\end{longtable}
\paragraph{\code{src/hamiltonians.py}}
\lstinputlisting{../src/hamiltonians.py}
\begin{longtable}{p{.12\textwidth}p{.79\textwidth}}
\toprule Lines & Mathematical and implementation commentary\\\midrule\endhead
1--7 & Use COO triplets and CSR storage rather than a dense $d\times d$ allocation. QuTiP wraps the final sparse operator for evolution.\\
8--19 & Store basis and physical parameters; require finite $L>0$, $\omega_0\geq0$ and finite TLS position. Build the free term and interaction separately.\\
20--27 & For each occupation compute $\sum_m|k_m|n_m$, then place it on ground and add $\omega_0$ on excited TLS. The ground-state energy origin is explicit, \eqref{eq:entries}.\\
28--35 & Select $f_m=\sqrt{|k_m|}$ or $1$; reject other labels. The signed momentum survives in the spatial phase even though the profile uses $|k|$.\\
36--37 & Compute $u_m=\ii f_m e^{-\ii k_mx_{\mathrm{tls}}}/\sqrt L$, \eqref{eq:H}. $D$ is not multiplied yet, allowing the same interaction matrix to be scaled.\\
38--38 & Allocate row/column/value lists and a mode tag per transition. Tags support exact interaction-energy attribution in \eqref{eq:energytriplets}.\\
39--42 & Iterate occupations, TLS labels and changed modes. The matrix column is the source $2i+s$; the target row flips the TLS to $1-s$.\\
43--48 & Annihilation requires positive occupation. Under RWA it also requires source TLS $g$. All retained bases are downward closed, so dictionary lookup of the decremented tuple is valid.\\
49--50 & Record $u_m\sqrt{n_m}$ and the changed-mode tag. The factor is the bosonic matrix element, with the imposed complex annihilation phase.\\
51--54 & Creation is allowed for either TLS without RWA and only source $e$ under RWA. The incremented tuple is tested, not assumed present.\\
55--59 & Omit targets outside the chosen space, exactly implementing projection. Retained creation receives $u_m^*\sqrt{n_m+1}$, conjugate to its reverse annihilation.\\
60--65 & Keep numerical triplet/tag arrays for diagnostics and convert COO to CSR. Zero-valued triplets can remain recorded, notably for sqrt at $k=0$.\\
66--70 & Require finite real $D$ and return $H_0+DV$ as a QuTiP operator. Hermiticity follows from paired transitions and real scaling, not solver output normalization.\\
\bottomrule\end{longtable}
\paragraph{\code{src/fidelities.py}}
\lstinputlisting{../src/fidelities.py}
\begin{longtable}{p{.12\textwidth}p{.79\textwidth}}
\toprule Lines & Mathematical and implementation commentary\\\midrule\endhead
1--6 & This module implements exactly two squared metrics; ambient full-space vectors and a ratio field proxy are not constructed.\\
7--16 & Accept a QuTiP ket or one-dimensional array of the basis dimension. Normalize the entire nonzero finite ket for metric evaluation, retaining raw norm diagnostics elsewhere.\\
17--21 & Reshape alternating components to $C$ with shape $(B,2)$, then compute $C^{\mathsf T}C^*$ as in \eqref{eq:rho}. For a raw ket its trace is the squared norm.\\
22--26 & Start the union with the first physical momentum list and its identity map. Internal union order need not be sorted because occupation keys are canonicalized.\\
27--38 & Match each second-grid mode with absolute tolerance $10^{-12}$ and zero relative tolerance. Reuse its union index, append if absent, or reject an ambiguous match.\\
39--46 & For each tuple, discard zero occupations and sort positive $(\text{union index},n)$ pairs. Store the two TLS amplitudes; missing modes are vacuum implicitly and vacuum has an empty key.\\
47--52 & Normalize both kets, align physical modes, and build occupation dictionaries. This low-level API cannot check box length or Hamiltonian conventions, which remain caller responsibilities.\\
53--54 & Sum the conjugate first-state times second-state TLS dot product over every shared physical key, \eqref{eq:fstate}. Do not normalize on that intersection; example \eqref{eq:intersectionexample} would otherwise be wrong.\\
55--59 & Reduce the individually normalized kets, square QuTiP root fidelity, and return the two metrics clipped for numerical overshoot. Vacuum embedding leaves these atomic reductions unchanged.\\
60--67 & Require both histories, exactly matching output arrays and equal $L$ to absolute $10^{-12}$. These checks precede pairing, preventing silent comparison of different time samples or boxes.\\
68--72 & Compare matching stored states and return two one-dimensional arrays of length $N_t$. The simulated comparator is an independent input, not the ambient embedding choice.\\
\bottomrule\end{longtable}
\paragraph{\code{src/energy_profile.py}}
\lstinputlisting{../src/energy_profile.py}
\begin{longtable}{p{.12\textwidth}p{.79\textwidth}}
\toprule Lines & Mathematical and implementation commentary\\\midrule\endhead
1--5 & Energy partitions consume the already-built Hamiltonian and propagated vector; there is no second interaction implementation.\\
6--11 & Cache the occupation matrix $n_{i,m}$ and photon-total axis $0,\ldots,\nu_{\max}$. The function $\nu(q)$ and projectors $Q_\nu$ are defined in \eqref{eq:energy-sector-projectors}; full can reach $MN$.\\
12--14 & Reshape the raw vector to $C_{i,s}$ of shape $(B,2)$ and form $|C_{i,s}|^2$. Calculations use raw quadratic expectations; see \eqref{eq:energy-raw-normalized} for normalization.\\
15--16 & The transpose occupation matrix times summed TLS probabilities gives $\langle N_m\rangle=\sum_{i,s}n_{i,m}|C_{i,s}|^2$. Multiply by $\omega_m=|k_m|$ to obtain the bare field term $F_m$ in \eqref{eq:energy-contributions}.\\
17--19 & For each record $\ell$, index source $c_\ell$, target $r_\ell$ and complex value $z_\ell$, then compute $\Re[v_{r_\ell}^*Dz_\ell v_{c_\ell}]$. The record and derivation are \eqref{eq:energy-record-definition}--\eqref{eq:energytriplets}; values contain $u_m$ but not $D$.\\
20--21 & Weighted \code{bincount} by tag $m_\ell$ sums directed terms into $W_m=\langle DV_m\rangle$. Both reverse orientations are already present; do not multiply this accumulated result by two, \eqref{eq:energy-reverse-pair}.\\
22--24 & Return $E_m=F_m+W_m$ and the separate bare TLS term $E_a=\omega_0\sum_i|C_{i,1}|^2$. Interaction is assigned to the mode columns by convention, not included in the TLS marker, \eqref{eq:energy-allocation}.\\
25--26 & Use the same experiment and Hamiltonian; the result is indexed by photon total, keeping both atomic components. It is independent of the mode-tag partition.\\
27--28 & Form the sparse matrix action and its row terms $\Re[v_q^*(Hv)_q]$, which define \eqref{eq:energy-sector-row-definition}. These terms include cross-sector coherences rather than projecting the Hamiltonian to its diagonal blocks.\\
29--30 & First sum the two TLS rows of each photon configuration, then group by total photon number. This gives $E_\nu=\langle(Q_\nu H+H Q_\nu)/2\rangle$, \eqref{eq:energysectors}; it is not conditional energy \eqref{eq:energy-conditional}.\\
\bottomrule\end{longtable}
\paragraph{\code{src/experiment.py}}
\lstinputlisting{../src/experiment.py}
\begin{longtable}{p{.12\textwidth}p{.79\textwidth}}
\toprule Lines & Mathematical and implementation commentary\\\midrule\endhead
1--15 & One common engine imports grid, basis, Hamiltonian and diagnostics. \code{momentum_modes} remains importable from this module for API convenience.\\
16--30 & Validate finite positive $T,\code{dt}$. Generate nominal output spacing with a roundoff tolerance, always including zero and exact $T$; if $\code{dt}>T$, return both endpoints.\\
31--42 & Build configuration from the dictionary API and call the same grid resolver as propagation. Optional selection requires packet information; dimension is evaluated without enumerating a basis.\\
43--50 & Use exact dimensions \eqref{eq:dimensions} with the selected count. Reject unknown basis names; no dense allocation is needed to estimate exponential full-space growth.\\
51--55 & Count outputs and estimate $16d$ bytes per ket, retained history depending on storage, two retained vector copies and the sparse-entry bound, \eqref{eq:memory}.\\
56--79 & Return base/selected representations and heuristic feasibility. Printed output exposes other-control implications and explicitly names excluded memory costs; it does not enforce a solver allocation limit.\\
80--87 & Forward all grid, cap, selection and storage inputs from the dataclass to the dictionary estimator without a separate estimate formula.\\
88--104 & Copy physical dictionaries, resolve base/selected grids, generate output times, build basis and projected Hamiltonian, then prepare the ket. Coefficient arrays are absent until propagation.\\
105--116 & Run \code{sesolve} with BDF/default tolerances, raw output normalization and final-state storage. Build a contiguous vector array; alternating ground/excited slices are views, shape $(N_t,B)$ or $(B,)$ for final-only.\\
117--121 & Find tuples with total exactly one, then identify their occupied mode. This works across all three basis orderings and gives an empty list at $N=0$.\\
122--138 & Require propagation, promote final-only arrays internally, and compute raw norm, TLS excitation, signed $1g$ populations and all photon totals. Sum identities are \eqref{eq:populationchecks}.\\
139--146 & The first argument is an output index, not a physical time. Extract only $1g$ amplitudes and multiply by positive Fourier phases and $L^{-1/2}$, \eqref{eq:wavefunction}.\\
147--160 & Resolve diagnostic requests: all retained states for None, final for $-1$, or nearest stored output for an in-range time. No interpolation occurs; final-only runs reject earlier requests.\\
161--164 & Reduce each requested raw ket. None returns a list, an explicit time one matrix; final-only with None returns a one-element list.\\
165--168 & Return the real excited diagonal of the raw TLS reduction. None returns an array, an explicit time a scalar; norm drift remains visible.\\
169--173 & Normalize the atomic reduction by the full squared ket norm before natural-log von Neumann entropy. This is \eqref{eq:entropy}, bounded by $\log2$ for normalized valid states.\\
174--178 & Apply the same sparse $H_0+DV$ and take the real raw expectation. Return a time array or scalar without introducing a second energy-origin convention.\\
179--184 & For requested kets, return $(k,E_m,0.0,E_a)$ with histories or one profile. The zero is a retained tuple-interface placeholder, not a supplemental atomic mode.\\
185--188 & Return the photon-total axis and row-grouped energy partition. The time selector follows the same rules as other diagnostics; summed energy agrees with direct expectation.\\
\bottomrule\end{longtable}
\paragraph{\code{experiment/_common.py}}
\lstinputlisting{../experiment/_common.py}
\begin{longtable}{p{.12\textwidth}p{.79\textwidth}}
\toprule Lines & Mathematical and implementation commentary\\\midrule\endhead
14--18 & \code{run}: Construct one \code{Experiment}, propagate, then compute observables in that order. Return the numerical object so notebook analyses can reuse its kets and Hamiltonian.\\
21--25 & \code{save_arrays}: Create the output directory and write NPZ arrays plus \code{configuration_json} from the base dataclass. Non-JSON values stringify; ordinary numerical arrays can be loaded without pickle.\\
28--53 & \code{parser}: Shared CLI names use $D,L,M,k_0,\sigma_k,x_0,N,T$. Default profile is sqrt. Explicit $M$ requires its boolean flag; default CLI grid control uses cutoffs.\\
56--65 & \code{config_from_args}: Construct the three dictionaries and dataclass with the current conventions. Both count and cutoff inputs may be stored, but only the selected control is active; historical values are not automatically converted.\\
68--77 & \code{finish}: With an output path, save the figure and same-stem NPZ, then close it. Otherwise display the figure. Backend selection remains the environment or notebook responsibility.\\
\bottomrule\end{longtable}
\paragraph{\code{experiment/atom_evolution.py}}
\lstinputlisting{../experiment/atom_evolution.py}
\begin{longtable}{p{.12\textwidth}p{.79\textwidth}}
\toprule Lines & Mathematical and implementation commentary\\\midrule\endhead
8--10 & \code{run_atom_evolution}: Replace the configuration by each scheme and force history storage. Return a dictionary of experiments, so excitation, entropy and notebook fidelity derive from the same stored trajectories.\\
13--24 & \code{plot_atom_evolution}: Plot raw $P_e$ and normalized TLS entropy against output times for every returned scheme. This plot does not itself propagate or compute cross-state fidelity.\\
27--37 & \code{main}: Parse current CLI inputs, report relevant estimates before simulation, call the module numerical routine, then plot/export. The main guard preserves importability from notebooks.\\
\bottomrule\end{longtable}
\paragraph{\code{experiment/compare_mode_selection.py}}
\lstinputlisting{../experiment/compare_mode_selection.py}
\begin{longtable}{p{.12\textwidth}p{.79\textwidth}}
\toprule Lines & Mathematical and implementation commentary\\\midrule\endhead
11--44 & \code{run_mode_selection}: Use an unselected total-cap comparator at each coupling. Sweep schemes and both selection flags; return $(S,2,Q)$ means/initial metrics. At one fixed $D$, scan windows to return $(S,P,A)$ mean-only heatmaps.\\
47--67 & \code{plot_mode_selection}: Plot selection on/off sweeps and window maps for both metrics. Transpose the stored $(P,A)$ slice for plotting so the horizontal axis remains photon window and vertical axis atom window.\\
70--83 & \code{main}: Parse current CLI inputs, report relevant estimates before simulation, call the module numerical routine, then plot/export. The main guard preserves importability from notebooks.\\
\bottomrule\end{longtable}
\paragraph{\code{experiment/energy_profile.py}}
\lstinputlisting{../experiment/energy_profile.py}
\begin{longtable}{p{.12\textwidth}p{.79\textwidth}}
\toprule Lines & Mathematical and implementation commentary\\\midrule\endhead
8--17 & \code{run_energy_profile}: Require history, propagate once, and attach mode/TLS, photon-number and total energies to the experiment. Their shapes and sum identities are documented above.\\
20--34 & \code{plot_energy_profile}: Plot time--momentum and time--photon-total heatmaps plus total/TLS energy. Negative partition values are energies, not invalid probabilities.\\
37--68 & \code{animate_energy_profile}: Reuse existing arrays for animation and optional GIF output, with stride and frame rate affecting display only. The marker at momentum zero represents separate TLS energy, not another oscillator.\\
71--81 & \code{main}: Parse current CLI inputs, report relevant estimates before simulation, call the module numerical routine, then plot/export. The main guard preserves importability from notebooks.\\
\bottomrule\end{longtable}
\paragraph{\code{experiment/full_cumulative_fidelity.py}}
\lstinputlisting{../experiment/full_cumulative_fidelity.py}
\begin{longtable}{p{.12\textwidth}p{.79\textwidth}}
\toprule Lines & Mathematical and implementation commentary\\\midrule\endhead
11--25 & \code{run_cap_convergence}: Require at least two strictly increasing caps; run total-cap final-only trajectories at each. Compare adjacent list entries initially and finally using physical embedding; retain all run objects.\\
28--37 & \code{plot_cap_convergence}: Two metric panels distinguish initial and final fidelities. The horizontal value is the lower cap in each adjacent pair, not necessarily a comparison of $N$ with $N+1$.\\
40--51 & \code{main}: Parse current CLI inputs, report relevant estimates before simulation, call the module numerical routine, then plot/export. The main guard preserves importability from notebooks.\\
\bottomrule\end{longtable}
\paragraph{\code{experiment/scattering.py}}
\lstinputlisting{../experiment/scattering.py}
\begin{longtable}{p{.12\textwidth}p{.79\textwidth}}
\toprule Lines & Mathematical and implementation commentary\\\midrule\endhead
12--26 & \code{run_scattering}: Preserve the user dictionary API and append grid-control inputs. Run shared dynamics; optional CSV contains observable columns only and overwrites its fixed repository results path.\\
29--47 & \code{plot_scattering}: Plot initial/middle/final $1g$ Fourier densities and directional/zero/TLS/norm populations. The interaction marker is $-x_{\mathrm{tls}}$, derived in \eqref{eq:position}.\\
50--59 & \code{main}: Parse current CLI inputs, report relevant estimates before simulation, call the module numerical routine, then plot/export. The main guard preserves importability from notebooks.\\
\bottomrule\end{longtable}
\paragraph{\code{experiment/sweep_D_fidelity.py}}
\lstinputlisting{../experiment/sweep_D_fidelity.py}
\begin{longtable}{p{.12\textwidth}p{.79\textwidth}}
\toprule Lines & Mathematical and implementation commentary\\\midrule\endhead
11--29 & \code{run_coupling_sweep}: Validate a finite nonempty one-dimensional $D$ array and force histories. At each $D$, propagate the chosen comparator and each candidate, reusing it for an identical scheme. Return $(S,Q)$ sample means and initial metrics, \eqref{eq:mean}.\\
32--42 & \code{plot_coupling_sweep}: Plot sample means versus physical $D$ and dashed initial values, separately for both metrics. The unsuffixed result is neither final fidelity nor a time-integral quadrature.\\
45--57 & \code{main}: Parse current CLI inputs, report relevant estimates before simulation, call the module numerical routine, then plot/export. The main guard preserves importability from notebooks.\\
\bottomrule\end{longtable}


\section{Comparison with the student implementation and validation}
\subsection{Exact convention conversions and finite-model changes}
\paragraph{Coupling, field phase and state conversion.}
An exact comparison first needs identical physical modes, bases and dispersion.
Call the historical parameter named \code{g} $\gamma_U$; its annihilation
coefficient was $\gamma_U\sqrt{|k_m|}e^{-\ii k_mx_a}$.
For sqrt and equal position parameters, matching magnitudes requires
\begin{equation}
 D=\sqrt L\,\gamma_U.
 \label{eq:Dconversion}
\end{equation}
The box factor means equal numerical old/new couplings need not agree.
The factor $\ii$ is a field phase convention implemented by
\begin{equation}
 U=e^{-\ii\pi\mathcal N_f/2},\quad Ua_mU^\dagger=\ii a_m,\quad
 Ua_m^\dagger U^\dagger=-\ii a_m^\dagger,\quad H_D=UH_UU^\dagger.
 \label{eq:gauge}
\end{equation}
Since $U\ket{\bm n;s}=(-\ii)^\nu\ket{\bm n;s}$, applying a ladder before
and after the diagonal phase proves the rules. All occupation projectors
commute with $U$, so the equality survives truncation. Matched trajectories
also require
\begin{equation}
 \ket{\psi_D(0)}=U\ket{\psi_U(0)},\qquad
 \ket{\psi_D(t)}=U\ket{\psi_U(t)}.
 \label{eq:trajectorygauge}
\end{equation}
Fixed-photon-number preparation changes only by a global phase. Coherent
preparation instead requires
\begin{equation}
 \alpha_D=-\ii\alpha_U,
 \label{eq:alphaconversion}
\end{equation}
including after projection. Keeping the same complex $\alpha$ while changing
the Hamiltonian phase is not a matched-state conversion. Correctly related
trajectories have equal TLS reductions and number populations; phase-sensitive
field observables must be transformed too. Fidelities are invariant when the
same $U$ acts on both compared states. Transforming only one member can
produce a convention-dependent overlap.

A displaced TLS additionally satisfies
\begin{equation}
 H(x)=U_xH(0)U_x^\dagger,\quad
 U_x=\exp\!\left(\ii x\sum_m k_mn_m\right),\quad
 U_xa_mU_x^\dagger=e^{-\ii k_mx}a_m.
 \label{eq:translationgauge}
\end{equation}
This checks both transition phases simultaneously. It does not alter the
Fourier-coordinate interpretation in equation~\eqref{eq:position}.

\paragraph{Packet-width and selection conversion.}
The historical amplitude was $\exp[-\sigma_U^2(k-k_0)^2/2]$.
Equating exponents with equation~\eqref{eq:packet} gives
\begin{equation}
 \sigma_k=\frac1{\sqrt2\sigma_U},\qquad
 x_0=x_{\mathrm{photon},U},\quad k_0=k_{\mathrm{photon},U}.
 \label{eq:width}
\end{equation}
The old width multiplies momentum differences; the new width is a momentum
standard-deviation parameter. Increasing old $\sigma_U$ narrows the momentum
packet. Normalization on identical modes then gives the same coefficients.

Historical selection radii were $w_U\sigma_U$, whereas new radii are
$w_D\sigma_k$. Equal numerical radii additionally require
\begin{equation}
 w_D=\frac{w_U\sigma_U}{\sigma_k}=\sqrt2\,w_U\sigma_U^2.
 \label{eq:windowconversion}
\end{equation}
This is a conversion in the historical numerical units. Envelope conversion
alone does not match a selected calculation: the window radii and nearest-mode
fallback must also be checked. Current factors mean multiples of the actual
momentum width, rather than automatically reproducing old factor values.

\paragraph{Implicit UV regulation and the changed grid.}
The historical folded grid used $L$ and requested $M_U$. For even
$M_U=2J$, its retained indices after zero deletion were
\begin{equation}
 \{-J+1,\ldots,-1,1,\ldots,J\},\quad |\K_U|=M_U-1,
 \qquad k_{\max}=J\Delta k=\frac{\pi M_U}{L}.
 \label{eq:oldgrid}
\end{equation}
The positive endpoint was unpaired. Full bandwidth is approximately
$2\pi M_U/L$ and radial UV approximately $\pi M_U/L$: these are different
uses of ``cutoff''. Current UV means largest retained $|k|$.
At $M_U=16,L=20$, old indices were $-7,\ldots,-1,1,\ldots,8$.
New explicit $M=15$ gives $-7,\ldots,0,\ldots,7$; $M=17$ gives
$-8,\ldots,0,\ldots,8$. Neither reproduces the old set. The symmetric
cutoff branch also cannot reproduce its unpaired endpoint.

This is an intentional change of finite model, not a stylistic adjustment
or a coupling redefinition. The new UV values at $L=20$ are $7\pi/10$ and
$4\pi/5$ respectively. If one wants a fixed physical UV while increasing
$L$ and refining momentum spacing, cutoff control states that intention
explicitly and lets the count follow. Positive IR excludes zero only by an
explicit band choice, not a hard-coded special case.

\paragraph{Parameter migration and compatibility boundaries.}
The numerical conversion rules can be located directly in the current inputs:
\begin{longtable}{p{.25\textwidth}p{.27\textwidth}p{.39\textwidth}}
\toprule Historical input & Current input & Meaning and required check\\\midrule\endhead
\code{length} & Atomic \code{L} & Same physical length in natural units;
the new momentum list must still be checked explicitly.\\
\code{modes} & Config \code{M} or cutoffs & Odd actual count under explicit
control; it is not the historical pre-deletion counter.\\
\code{g} & Atomic \code{D} & $D=\sqrt L\,\gamma_U$ for sqrt and matched
units; coherent phase also requires~\eqref{eq:alphaconversion}.\\
\code{w_atom}, \code{x_atom} & Atomic \code{omega_0}, \code{x_tls} & Same
frequency and Hamiltonian phase parameter. Do not inherit the old nonzero
default position when intending a TLS at zero.\\
\code{k_photon}, \code{x_photon} & Photon \code{k_0}, \code{x_0} & Same
signed packet centre and phase parameter at matched modes.\\
\code{sigma_photon} & Photon \code{sigma_k} & Reciprocal-width conversion
\eqref{eq:width}, not a direct rename.\\
\code{state} class & Photon \code{state} string & \code{number} or
\code{coherent}; explicit \code{n} or \code{alpha} in the photon dictionary.\\
\code{atom_state} & Atomic \code{initial_state} & Current preparation
interface accepts $g,e,+,-$ labels. The former arbitrary two-component
sequence is not an accepted label in this interface.\\
\code{excitation_cap} & Config \code{n_max} & Cap meaning depends on the
basis: per-mode in full, total in total-cap, ray occupation in truncated.\\
\code{t}, \code{dt} & Time \code{T}, \code{dt} & Same duration and nominal
sampling intention, but a changed output-time construction.\\
\code{hbar}, \code{c} & Fixed natural units & The current API assumes both
equal one; it does not expose arbitrary dimensional constants.\\
\bottomrule
\end{longtable}
The old special single-mode branch placed a mode at the packet momentum.
Current explicit $M=1$ instead means the zero mode, as specified by the
symmetric-grid rule. Legacy classes such as \code{Config} and the separate
Hamiltonian/state subclasses are replaced by the shared configuration and
engine. Experiments use the canonical modules in \code{experiment/}; the historical
internal class API and argument syntax are not preserved as aliases. The user's
dictionary-based scattering interface is retained.

Sampling changes also have a concrete consequence. At $T=4$ and
\code{dt=0.05}, the old construction requested $80$ outputs with spacing
$4/79$, whereas the current construction requests $81$ outputs with spacing
$0.05$. A mean over those samples is a different discrete statistic even
if both solvers approximate the same continuous trajectory. In the old
construction, \code{dt} was used to choose a count rather than enforced as
an output spacing.

\paragraph{Comparative summary and corrected overlap.}
\begin{longtable}{p{.22\textwidth}p{.31\textwidth}p{.38\textwidth}}
\toprule Item & Student implementation & Patch and consequence\\\midrule\endhead
Grid & Requested count reduced by zero deletion; unpaired endpoint
possible. & Exact odd count with zero or an explicit radial band; reports
resulting count and effective cutoffs.\\
Coupling & Historical parameter, real coefficient at zero position. &
$\ii Df_m/\sqrt L$ with signed phase. Exact conversion at fixed grid is
\eqref{eq:Dconversion}--\eqref{eq:trajectorygauge}. Uri notebooks use sqrt.\\
Packet & Multiplicative width in exponent. & Momentum width in denominator,
conversion~\eqref{eq:width}; all selected momentum signs enter preparation.\\
Assembly & Search among pairs of states. & Dictionary targets and bosonic
transitions, $O(Md)$ assembly; shared entries fixed by~\eqref{eq:entries}.\\
Indexing & Separate classes and scheme-specific indices. & Occupation tuples,
alternating TLS entries and comparisons of physical configurations.\\
Cross-basis overlap & Different basis types compared through one-occupied-mode
configurations. & All shared occupations in implicit full, physical momentum
alignment, and vacuum embedding of absent modes.\\
Metrics & Global/atomic metrics plus a ratio used as a field proxy. & Only
squared state/atomic fidelities; initial differences separately visible.\\
Times & \code{linspace(0,T,int(T/dt))} and output renormalization. & Nominal
\code{dt}, exact final time, raw norms and explicit tolerances. Sample means
may change even with similar trajectories.\\
Control & Heterogeneous scripts and inputs. & Callable modules, five notebooks
and common CLI. Former \code{codex/} migrated and removed.\\
\bottomrule
\end{longtable}
A minimal counterexample to the old cross-type overlap is
$\ket{1,1;g}$ in full and total-cap with $N=2$. The physical overlap is one;
a sum over one occupied mode gives zero. The corrected pair returns both
fidelities equal to one. This fixes a comparison, not a propagated trajectory.

At identical grids, number preparations are identical in every basis that
contains them, and remain identical for $D=0$. Coherent projections can
have equation~\eqref{eq:initialf} below one. Since occupation projectors
commute with $H_0$, their mutual overlap remains at that initial value during
free evolution on the shared grid. This diagnoses preparation differences
without mistaking them for interaction-induced error.

Reproducing old figures requires the exact mode list and cap/basis, coupling
magnitude and phase, envelope and coherent phase, selection radii, output
sampling and statistic. Parameter labels alone cannot establish reproduction
when the finite grids or initial projections differ.

\subsection{Mathematical verification and limits of the conclusions}
\paragraph{Independent finite-dimensional oracles.}
Tests construct small oscillator/TLS tensor operators independently of the
transition loops, select rows/columns for each occupation basis and check
\begin{equation}
 H_{\mathrm{code}}=P_\B H_{\mathrm{tensor}}P_\B,
 \qquad H_{\mathrm{code}}=H_{\mathrm{code}}^\dagger.
 \label{eq:tensorcheck}
\end{equation}
Checks cover three bases, two profiles, RWA on/off and zero/nonzero positions.
They also test parity, RWA excitation conservation and the simultaneous
annihilation/creation phase identity~\eqref{eq:translationgauge}.

An analytic check takes $L=2\pi$, modes $\{-1,0,1\}$, sqrt,
$\omega_0=1$, RWA, $N\geq1$ and input $\ket{0;e}$. At zero TLS position,
\begin{equation}
 \ket B=(\ket{1_-;g}+\ket{1_+;g})/\sqrt2,\qquad
 \ket{\mathrm{dark}}=(\ket{1_-;g}-\ket{1_+;g})/\sqrt2.
\end{equation}
The dark and zero-mode states decouple. In $(\ket{0;e},\ket B)$,
\begin{equation}
 H=I+\begin{pmatrix}0&\ii G\\-\ii G&0\end{pmatrix},\qquad
 G=\frac{\sqrt2D}{\sqrt L},\qquad P_e(t)=\cos^2(Gt).
 \label{eq:rabioracle}
\end{equation}
This checks the two-mode factor, box normalization and evolution in all bases.
Independent $e^{-\ii Ht}\psi_0$ comparisons check BDF and final-only storage.

Fidelities are tested against explicit small common-space embeddings with
reordered grids, selected subsets, caps and shared multimode configurations.
Atomic fidelity is checked against equation~\eqref{eq:qubit}. The $(1,1)$
case protects the correction; orthogonal field directions with identical TLS
protect the distinction between metrics. Photon populations sum to the
ket norm, and the total Hamiltonian expectation is checked for conservation. Gaussian, vacuum,
coherent and zero-mode preparations are tested separately from evolution.

Integration tests exercise all five experiment module CLIs, their exports,
and the callable functions. The five notebooks are checked for valid Python
cells and the requested grid/coupling conventions. Tabular diagnostics are
checked against the raw arrays, including final-only labels, absent grid
equivalences and preservation of norm drift and numerical precision. A documentation rebuild
does not replace rerunning the suite after numerical code changes. In an installed development environment, use:
\begin{lstlisting}[language={}]
python -m pytest tests -q
\end{lstlisting}
Documentation compiles from its sources and regenerates line comments from
current files. Derivations added here explain the current model; they do
not introduce additional solver features.

\paragraph{Integration error versus projection error.}
Norm and energy conservation are necessary checks, but two different unitary
evolutions can preserve both. Tighter tolerances or an independent exponential
reference on small systems test integration error more directly.
For an interpolated numerical trajectory with
$r(t)=\ii\dot v_{\mathrm{num}}-H_{\B}v_{\mathrm{num}}$, variation of constants gives
\begin{equation}
 \|v_{\mathrm{num}}(t)-v(t)\|
 \leq\|v_{\mathrm{num}}(0)-v(0)\|+\int_0^t\|r(s)\|\,ds.
 \label{eq:numericalbound}
\end{equation}
The engine reports sampled diagnostics, not an evaluated residual integral;
this equation specifies what a residual-based certificate would require.

For projection, take a larger finite space with Hamiltonian $H$, smaller
projector $P$, and $Q=I-P$. Let $\psi$ solve the larger evolution and $\psi_P$
solve $PHP$. For $e=P\psi-\psi_P$,
\begin{equation}
 \ii\dot e=PHP\,e+PHQ\psi,\qquad
 \|e(t)\|\leq\|e(0)\|+\int_0^t\|PHQ\psi(s)\|\,ds.
 \label{eq:projectionbound}
\end{equation}
The driving term depends on omitted dynamics, not just final boundary
population. Normalized coherent projections can give $e(0)\ne0$.
Cap comparisons are useful, but a small final boundary population alone
is not a rigorous dynamical error certificate. The modules do not evaluate
this bound.

\paragraph{Regulator convergence and finite-box scattering.}
The checks establish finite-model consistency, not physical continuum convergence.
A study still needs separate variations of
\begin{equation}
 N\to N',\quad k_{\mathrm{UV}}\to k'_{\mathrm{UV}},\quad
 k_{\mathrm{IR}}\to k'_{\mathrm{IR}},\quad L\to L',\quad
 (\mathrm{rtol},\mathrm{atol})\to(\mathrm{rtol}',\mathrm{atol}').
 \label{eq:regulators}
\end{equation}
Truncated omits multimode configurations at every $N$; its large-cap limit
is not full Fock space. Conversely, RWA number preparations can occupy a
small invariant excitation sector in which larger caps make no difference.
That agreement does not test arbitrary high-photon dynamics.

Changing $L$ at fixed explicit count changes spacing, UV, resonant-mode
availability and $D/\sqrt L$ together. Holding physical $D$ and cutoffs fixed
while changing $L$ is a different discretization study. For a finite smooth
band the associated Riemann sum is
\begin{equation}
 \frac1L\sum_{k_m\in\K}f_m^2\longrightarrow
 \frac1{2\pi}\int_{k_{\mathrm{IR}}\leq|k|\leq k_{\mathrm{UV}}}f(k)^2\,dk.
 \label{eq:continuumweight}
\end{equation}
This explains why per-mode $L^{-1/2}$ normalization and a growing mode count
belong together. It does not prove convergence as UV increases without bound.

Free phases $e^{-\ii|j|\Delta k t}$ recur after $t=L$. Interactions can change
exact recurrence times, but a periodic packet can return to the interaction
region after circulating around the box. Scattering observations should
separate incoming/outgoing packets before such returns and be checked in
larger boxes. Increasing $T$ without this check can turn single-passage
scattering into repeated interaction. Initial negative/zero components and
non-RWA multiphoton sectors must also be accounted for before assigning
asymptotic coefficients.

Historical plots become numerical references only after their actual grids,
preparations, coupling conversions and statistics have been reconstructed.
The equations, visible notebook inputs and configuration metadata make these
comparisons reproducible without treating a changed finite model as a
code-style refinement.
\end{document}
